% coha_wall_crossing_platonic.tex
% Platonic-ideal reconstitution of Vol III Kontsevich-Soibelman wall-crossing
% content. Written in the Russian-school style (Kontsevich, Soibelman,
% Schiffmann, Vasserot, Davison). Strictest honest scope, AP-CY5 / AP-CY7 /
% AP-CY8 compliant. No AI attribution; all work attributed to Raeez Lorgat.
%
% This chapter is the healing inscription following the 2026-04-16 adversarial
% attack on the CoHA wall-crossing claims in toric_cy3_coha.tex. It isolates
% the algebra-vs-coalgebra distinction, the positive-half-vs-full-Yangian
% distinction, and the motivic-Hall-Lie-algebra-vs-chiral-convolution-dgLA
% distinction that the earlier inscription slurs over in a few places.
%
% Chapter placement: sibling of toric_cy3_coha.tex in the CY landscape part
% (see main.tex \ref{part:cy-landscape}).

\chapter{Wall-Crossing and CoHA: the Algebra-vs-Coalgebra Ledger}
\label{ch:coha-wall-crossing-platonic}

\section*{Foreword: what this chapter is}

The preceding chapter \ref{ch:toric-coha} assembles a chiral quantum
group datum from the critical CoHA, the Maulik--Okounkov $R$-matrix, the
Drinfeld double, and the chart-gluing chiral algebra $A_X$. That assembly
is correct in its structural outline but passes quickly through three
distinctions that this chapter makes ledger-strict:

\begin{enumerate}[label=\textup{(\roman*)}]
  \item \textbf{Algebra versus coalgebra.}
    The cohomological Hall algebra $\mathcal{H}(Q, W)$ is a
    $\mathbb{Z}$-graded associative algebra \emph{without internal
    differential}. The bar complex $B^{\mathrm{ord}}(\mathcal{H})$ is a
    coassociative coalgebra constructed \emph{from} $\mathcal{H}$. The
    phrase ``$D^2 = 0$ in the CoHA differential'' in
    Theorem~\ref{thm:wall-crossing-mc} is shorthand for
    ``$D^2 = 0$ in the bar complex of $\mathcal{H}$'' or for
    ``$\partial_W^2 = 0$ in the Ginzburg dg algebra $\Pi(Q, W)$''; these
    are distinct objects and the shorthand suppresses the passage. This
    chapter separates them.
  \item \textbf{Positive half versus full Yangian.}
    The Schiffmann--Vasserot theorem identifies
    $\mathcal{H}(\mathbb{C}^3) \simeq Y^+(\widehat{\mathfrak{gl}}_1)$,
    the \emph{positive half}. The full affine Yangian
    $Y(\widehat{\mathfrak{gl}}_1)$ is the Drinfeld double of
    $\mathcal{H}(\mathbb{C}^3)$ via a non-degenerate Hopf pairing. For
    $X = \mathbb{C}^3$ this double is proved (Procha\'zka--Rap\v{c}\'ak);
    for general toric $X$ the double exists but the identification
    with a specific affine super Yangian is conditional on a Hopf-pairing
    lemma which we record as a proposition with precise hypotheses.
  \item \textbf{Motivic Hall Lie algebra versus chiral convolution dgLA.}
    Kontsevich--Soibelman wall-crossing lives in the motivic Hall Lie
    algebra $\mathfrak{g}_{\mathrm{KS}}(X)$ of the CY$_3$ category
    $D^b\mathrm{Coh}(X)$; its exponential is the motivic quantum torus
    and the wall-crossing automorphism is an exponential gauge action.
    The MC equation on the bar side of
    Theorem~\ref{thm:wall-crossing-mc} lives in the convolution dgLA
    $\mathfrak{g}^{\mathrm{mod}}_{A_X}$ of the chiral algebra $A_X$. The
    identification of these two dgLAs factors through the chiralization
    functor $\Phi$; it is proved for $\mathbb{C}^3$ (RSYZ plus
    the conifold bar-complex engine) and conditional on CY-A$_3$ outside
    that case.
\end{enumerate}

The structural yoga of the programme survives these sharpenings unchanged.
What changes is the precise hypothesis ledger, and it is that ledger which
this chapter lays down.

% ========================================================================
\section{Setup: CoHA, Ginzburg dg algebra, bar complex}
\label{sec:coha-setup-platonic}

Let $Q = (Q_0, Q_1, s, t)$ be a finite quiver with potential
$W \in \mathbb{C}Q_{\mathrm{cyc}}$ (cyclic path algebra). For toric
CY$_3$ geometries $X$ without compact $4$-cycles the McKay-correspondence
quiver $(Q_X, W_X)$ is the input.

\begin{definition}[Three distinct objects]
\label{def:three-coha-objects}
\ClaimStatusProvedHere
Associated to $(Q, W)$ we distinguish:
\begin{enumerate}[label=\textup{(\alph*)}]
  \item \textbf{Path algebra.} The path algebra $\mathbb{C}Q$ with
    relations $\partial W / \partial a_i = 0$ for each $a_i \in Q_1$
    (the Jacobi algebra
    $J(Q, W) = \mathbb{C}Q / (\partial W)$).
    This is an associative $\mathbb{Z}_{\geq 0}$-graded algebra
    (grading by path length), \emph{without} internal differential.
  \item \textbf{Ginzburg dg algebra.} The dg algebra $\Pi(Q, W)$
    with generators $Q_0 \cup Q_1 \cup Q_1^* \cup \{\epsilon_i\}_{i \in Q_0}$
    in degrees $0, 0, -1, -2$ and differential $\partial_W$
    determined by
    $\partial_W(a_i^*) = \partial W / \partial a_i$ and
    $\partial_W(\epsilon_i) = \sum_{a \in Q_1} [a, a^*]\, e_i$.
    This is a genuinely dg algebra: $\partial_W^2 = 0$ encodes the
    Jacobi identity for $W$ as a cyclic function.
  \item \textbf{Critical CoHA.} The cohomological Hall algebra
    \[
      \mathcal{H}(Q, W)
       \;=\;
      \bigoplus_{\mathbf{d} \in \mathbb{Z}_{\geq 0}^{Q_0}}
       \; H^{\bullet}_{G_\mathbf{d}}
        \bigl( \mathrm{Rep}(Q, \mathbf{d}),
         \phi_{W_\mathbf{d}} \bigr),
    \]
    a $\mathbb{Z}^{Q_0}$-graded associative algebra with Hall product
    (Kontsevich--Soibelman 2010, Davison--Meinhardt 2015). $\mathcal{H}$
    is an associative graded algebra: it carries the Hall product as
    $m_2$ and NO internal differential at the level of definition.
\end{enumerate}
\end{definition}

\begin{proof}[Justification]
(a) is literally the Jacobi algebra. (b) is Ginzburg 2006
\textup{arXiv:math/0612139}, Proposition~5.1.9; the identity
$\partial_W^2 = 0$ is a cyclic-derivation calculation (the double
derivative $\partial^2 W / \partial a_i \partial a_j$ is symmetric,
and $\partial_W^2$ contracts it against $[a_i^*, a_j^*] = 0$). (c)
is Kontsevich--Soibelman 2010 \textup{arXiv:1006.2706} Section~7,
specialized to the critical (with-potential) setting; cf.\
Davison--Meinhardt 2015 \textup{arXiv:1601.02479} for the
integrality theorem. Each object has its own grading and its own
role: the Jacobi algebra is the classical shadow, the Ginzburg dg
algebra is the dg model that carries the differential, the CoHA is
the cohomological invariant built \emph{from} the Ginzburg dg algebra
by taking Borel--Moore homology of the critical locus.
\end{proof}

\begin{remark}[Which object carries the differential]
\label{rem:which-carries-differential}
Among (a)--(c) of Definition~\ref{def:three-coha-objects}:
\begin{itemize}
  \item (a) has NO differential; it is a graded algebra.
  \item (b) HAS a differential $\partial_W$; it is a dg algebra. This is
    where $\partial_W^2 = 0$ is a nontrivial identity encoding the
    Jacobi algebra relations.
  \item (c) has NO internal differential; its grading is by the charge
    lattice $\mathbb{Z}^{Q_0}$, and its algebra structure is the
    cohomological Hall product built by taking $H^\bullet_{G_\mathbf{d}}$
    of vanishing cycles. The CoHA IS the output of a derived
    construction, but the output itself is a graded algebra, not a dg
    algebra.
\end{itemize}
Prose of the form ``$D^2 = 0$ in the CoHA differential'' is a
\emph{category mistake} unless the author has in mind either (i) the
Ginzburg dg algebra differential, or (ii) the bar differential on
$B^{\mathrm{ord}}(\mathcal{H})$. Passages that mix these silently
are the targets of Theorem~\ref{thm:coha-is-algebra-not-coalgebra}.
\end{remark}

\begin{definition}[Bar complex of the CoHA]
\label{def:bar-of-coha-platonic}
\ClaimStatusProvedHere
The ordered bar complex of the CoHA is
\[
  B^{\mathrm{ord}}(\mathcal{H})
   \;=\; T^c\bigl( s^{-1} \overline{\mathcal{H}} \bigr),
\]
where $\overline{\mathcal{H}} = \ker(\epsilon)$ is the augmentation
ideal of $\mathcal{H}$ (as a graded algebra over $\mathbb{C}$) and
$T^c$ is the tensor coalgebra with deconcatenation coproduct. The
bar differential $d_B = \sum_{k \geq 1} b_k$ encodes the Hall product
as $b_2$; since $\mathcal{H}$ is associative and carries no higher
$A_\infty$-operations at the level of definition, $b_k = 0$ for
$k \geq 3$ and $b_2$ is the dual of the Hall product.
\end{definition}

\begin{remark}[Four inputs, four outputs, one confusion]
\label{rem:four-inputs-four-outputs-platonic}
The four objects on the table are:
\begin{center}
\renewcommand{\arraystretch}{1.15}
\begin{tabular}{lll}
\toprule
Object & Structure & Carries $d^2 = 0$? \\
\midrule
Jacobi algebra $J(Q, W)$ & graded algebra & no (no differential) \\
Ginzburg dg alg $\Pi(Q, W)$ & dg algebra & yes (native) \\
Critical CoHA $\mathcal{H}(Q, W)$ & graded algebra & no (no differential) \\
$B^{\mathrm{ord}}(\mathcal{H})$ & dg coalgebra & yes (bar differential) \\
\bottomrule
\end{tabular}
\end{center}
The confusion of writing ``$D^2 = 0$ in the CoHA'' targets the wrong
object: the CoHA is a graded algebra, not a dg object. The bar complex
\emph{of} the CoHA is a dg coalgebra and does carry $D^2 = 0$. The
Ginzburg dg algebra is a dg algebra with $\partial_W^2 = 0$. The
three are linked by a chain of functors, not equated.
\end{remark}

% ========================================================================
\section{CoHA is an algebra, not a coalgebra}
\label{sec:coha-is-algebra}

\begin{theorem}[CoHA is an algebra, not a coalgebra]
\label{thm:coha-is-algebra-not-coalgebra}
\ClaimStatusProvedHere
Let $(Q, W)$ be a quiver with potential and
$\mathcal{H} = \mathcal{H}(Q, W)$ its critical CoHA. Then:
\begin{enumerate}[label=\textup{(\roman*)}]
  \item $\mathcal{H}$ is an associative $\mathbb{Z}^{Q_0}$-graded
    algebra over $\mathbb{C}$ with Hall product
    $m_H \colon \mathcal{H}_{\mathbf{d}_1} \otimes \mathcal{H}_{\mathbf{d}_2}
     \to \mathcal{H}_{\mathbf{d}_1 + \mathbf{d}_2}$;
  \item $\mathcal{H}$ carries NO internal differential in its definition
    as a graded algebra; the Borel--Moore homology $H^{\bullet}_{G_\mathbf{d}}$
    that builds $\mathcal{H}$ is taken \emph{after} passing to cohomology,
    not before;
  \item the bar complex $B^{\mathrm{ord}}(\mathcal{H}) = T^c(s^{-1} \bar{\mathcal{H}})$
    is a separate object: a coassociative dg coalgebra whose differential
    is the bar differential $d_B$, with $d_B^2 = 0$;
  \item the statement ``$D^2 = 0$ in the CoHA'' in the prose of
    Theorem~\textup{\ref{thm:wall-crossing-mc}} refers to the bar
    complex $B^{\mathrm{ord}}(\mathcal{H})$ or to the Ginzburg dg algebra
    $\Pi(Q, W)$, not to the CoHA itself.
\end{enumerate}
\end{theorem}

\begin{proof}
(i) is Kontsevich--Soibelman 2010 Theorem~1 or Schiffmann--Vasserot
2013 Theorem~A specialised to the critical setting. The Hall product
is the convolution $m_H = \pi_* \circ p^*$ on equivariant Borel--Moore
homology with vanishing-cycle coefficient, where $p$ is the
short-exact-sequence projection and $\pi$ is the quotient to the
target dimension vector. Associativity is Schiffmann--Vasserot 2013
Proposition~1.4.

(ii) is immediate from the construction: $\mathcal{H}_\mathbf{d}$
is defined as a cohomology group
$H^\bullet_{G_\mathbf{d}}(\mathrm{Rep}(Q, \mathbf{d}), \phi_{W_\mathbf{d}})$,
which is a graded vector space; vanishing cycles is applied once, then
equivariant cohomology is taken, producing a graded vector space
without residual differential. There is no chain-level dg structure
\emph{on} $\mathcal{H}$ itself; there is such a structure on the
Ginzburg dg algebra $\Pi(Q, W)$ whose Hochschild homology feeds in,
but $\mathcal{H}$ is the cohomological output.

(iii) is the standard bar-complex construction. Given an associative
graded algebra $A$, $B^{\mathrm{ord}}(A) = T^c(s^{-1} \bar A)$ is the
tensor coalgebra on the desuspended augmentation ideal, with
deconcatenation coproduct and differential $d_B = \sum_{k \geq 1} b_k$
where $b_2$ is dual to the algebra multiplication and $b_k = 0$ for
$k \geq 3$ when $A$ is strictly associative. $d_B^2 = 0$ is equivalent
to associativity of $A$ (Stasheff 1963, Loday--Vallette 2012 \S2.3).

(iv) by remaining consistent with (i)-(iii) we are forced to
re-interpret any chain-level identity $D^2 = 0$ in Theorem~%
\ref{thm:wall-crossing-mc} as living on the bar side or the Ginzburg
side, not on the CoHA. The Ginzburg-side interpretation is the one used
in the proof of Theorem~\ref{thm:wall-crossing-mc}(iii): the identity
$\partial_W^2 = 0$ in $\Pi(Q, W)$ is what encodes the Jacobi
$\partial W / \partial a_i = 0$ condition, and this is what enters the
MC equation via the dg Lie algebra associated to $\Pi(Q, W)$.
\end{proof}

\begin{remark}[Reconciliation with the prose of Theorem~\ref{thm:wall-crossing-mc}]
\label{rem:reconciliation-wall-crossing-prose}
Clause (iii) of Theorem~\ref{thm:wall-crossing-mc} reads:
\begin{quote}
  The Jacobi algebra condition $\partial W / \partial a_i = 0$ is
  equivalent to $D^2 = 0$ in the CoHA differential, which in turn is
  the MC equation $D\Theta + \frac{1}{2}[\Theta, \Theta] = 0$.
\end{quote}
The correct reading: the Jacobi algebra condition
$\partial W / \partial a_i = 0$ defines the zero locus of
$\partial_W$ in the Ginzburg dg algebra, equivalently
$\partial_W^2 = 0$ holds in $\Pi(Q, W)$. This translates to the MC
equation
$D\Theta_W + \frac{1}{2}[\Theta_W, \Theta_W] = 0$ in the dg Lie
algebra $(\Pi(Q, W))_{\mathrm{Lie}}$ associated to
$\Pi(Q, W)$ with MC element
$\Theta_W = \sum_i (\partial W / \partial a_i) \, a_i^*$ (Ginzburg
2006, Lemma~5.3.1). The CoHA enters only as the cohomological
invariant of this dg data; it is not itself a dg object, and the
MC equation is on $\Pi(Q, W)_{\mathrm{Lie}}$, not on
$\mathcal{H}(Q, W)$.
\end{remark}

% ========================================================================
\section{Chiralisation preserves algebra structure}
\label{sec:chiralisation-preserves-algebra}

\begin{theorem}[Chiralisation preserves the algebra-side structure]
\label{thm:coha-chiralization-preserves-algebra}
\ClaimStatusProvedHere
Let $X$ be a local CY$_3$ with quiver $(Q_X, W_X)$ and associated critical
CoHA $\mathcal{H}(Q_X, W_X)$. Let
$\Phi \colon D^b\mathrm{Coh}(X) \to \mathrm{ChirAlg}$ denote the
chiralisation functor of Theorem~CY-A$_3$ (proved in the
$\infty$-categorical framework; Theorem~\textup{\ref{thm:cy-to-chiral-d3}}).
Then the image $\Phi(X) = A_X$ is an $E_1$-chiral ALGEBRA, and the
CoHA $\mathcal{H}(Q_X, W_X)$ embeds as the positive-half subalgebra:
\[
  \mathcal{H}(Q_X, W_X)
   \;\hookrightarrow\;
  A_X
\]
as graded algebras, via the identification
$\mathcal{H}(Q_X, W_X) \simeq Y^+(\widehat{\mathfrak{g}}_{Q_X})$
(Schiffmann--Vasserot for $\mathbb{C}^3$, Rapcak--Soibelman--Yang--Zhao
for general toric without compact $4$-cycles) and the identification
of $Y^+$ with a subalgebra of $A_X$ generated by the positive simple roots.
In particular, chiralisation preserves ALGEBRA structure: an associative
graded algebra $\mathcal{H}$ lands in an $E_1$-chiral algebra $A_X$,
not in a chiral coalgebra.
\end{theorem}

\begin{proof}
The functor $\Phi$ is constructed in Theorem~\ref{thm:cy-to-chiral-d3}
as a symmetric-monoidal $\infty$-functor from $D^b\mathrm{Coh}(X)$
with Yoneda convolution to the $\infty$-category of $E_1$-chiral
algebras (for $d \geq 3$; $E_2$-chiral at $d \leq 2$, cf.\ FM43). Its
input is an associative monoidal structure (Yoneda convolution), and its
output is an $E_1$-algebra structure. There is no conversion to
coalgebra at any stage.

The embedding of $\mathcal{H}$ as the positive-half subalgebra of $A_X$
proceeds through the identification
$\mathcal{H}(Q_X, W_X) \simeq Y^+(\widehat{\mathfrak{g}}_{Q_X})$ (RSYZ;
Schiffmann--Vasserot for $\mathbb{C}^3$). The affine super Yangian
$Y(\widehat{\mathfrak{g}}_{Q_X}) = Y^+ \otimes Y^0 \otimes Y^-$ has its
positive Borel $Y^+$ sitting inside $A_X$ as the subalgebra generated by
positive-root mode operators; this inclusion is an inclusion of graded
associative algebras, preserved by $\Phi$ at the level of mode algebras.
No coalgebra structure enters: the chiralisation functor is algebra-to-%
algebra by construction.

The bar complex $B^{\mathrm{ord}}(A_X) = T^c(s^{-1} \bar A_X)$ is a
separate object, a dg coalgebra; this is the coalgebra resolution of
$A_X$ on the Koszul locus and should not be confused with $A_X$ itself.
The CoHA embeds in $A_X$; it does not embed in $B^{\mathrm{ord}}(A_X)$.
\end{proof}

\begin{remark}[What chiralisation does NOT do]
\label{rem:chiralisation-does-not}
The functor $\Phi$ does NOT produce a chiral coalgebra from a CY
category. It produces an $E_n$-chiral algebra ($n = 2$ at $d \leq 2$,
$n = 1$ at $d \geq 3$). The bar complex of that output is a chiral
coalgebra; but the output itself is an algebra. Any prose of the form
``$\Phi$ converts CoHA to a chiral coalgebra'' is a violation of
and.
\end{remark}

% ========================================================================
\section{KS wall-crossing as MC equation on the motivic Hall Lie algebra}
\label{sec:ks-mc-on-hall-dgla}

\begin{theorem}[KS wall-crossing is MC gauge on the motivic Hall dgLA]
\label{thm:ks-wall-crossing-mc-on-coha-dgla}
\ClaimStatusProvedHere
Let $X$ be a local CY$_3$ and $\mathfrak{g}_{\mathrm{KS}}(X)$ the
motivic Hall Lie algebra of $D^b\mathrm{Coh}(X)$ in the sense of
Kontsevich--Soibelman 2008 (Section~2.3 of \textup{arXiv:0811.2435}).
Then:
\begin{enumerate}[label=\textup{(\roman*)}]
  \item $\mathfrak{g}_{\mathrm{KS}}(X)$ is a nilpotent pro-finite
    dg Lie algebra with bracket determined by Euler forms of exact
    triangles; the exponential $\exp \colon \mathfrak{g}_{\mathrm{KS}}
    \to G_{\mathrm{KS}}$ is the motivic quantum torus.
  \item For generic stability parameters $\zeta, \zeta'$ in adjacent
    chambers separated by a wall $W$, the KS wall-crossing formula is the
    gauge transformation
    \[
      \Theta_{\zeta'}
       \;=\; \exp(\mathrm{ad}_{\alpha_W})
       (\Theta_\zeta)
      \qquad
      \text{in } \mathfrak{g}_{\mathrm{KS}}(X),
    \]
    where $\alpha_W \in \mathfrak{g}_{\mathrm{KS}}(X)$ is supported on
    the wall class $\gamma_W \in K_0(\mathrm{Coh}(X))$, and
    $\Theta_\zeta$ is the MC element whose coefficients are the
    motivic DT invariants of the $\zeta$-semistable objects.
  \item The MC equation
    $D\Theta_\zeta + \frac{1}{2}[\Theta_\zeta, \Theta_\zeta] = 0$
    in $\mathfrak{g}_{\mathrm{KS}}(X)$ is the closedness of the DT
    generating series against the motivic Hall product.
  \item The pentagon identity
    $\exp(\mathrm{ad}_{\alpha_5}) \circ \cdots \circ
     \exp(\mathrm{ad}_{\alpha_1}) = \mathrm{id}$
    is the $A_2$ cluster-mutation periodicity (five pentagon
    cycles) in $\mathfrak{g}_{\mathrm{KS}}(X)$.
\end{enumerate}
This MC equation lives on the ALGEBRA side: $\mathfrak{g}_{\mathrm{KS}}$
is a graded Lie algebra, the gauge action is algebra-preserving, and no
bar coalgebra enters at this level.
\end{theorem}

\begin{proof}
(i) is Kontsevich--Soibelman 2008 Section~2.3. The Hall Lie algebra
is freely generated by the $K$-theory classes of simple objects
with bracket $[e_\gamma, e_{\gamma'}] = (-1)^{\chi(\gamma, \gamma')}
(\chi(\gamma, \gamma') - \chi(\gamma', \gamma))\, e_{\gamma + \gamma'}$
(Euler form), nilpotent after completion in the Harder--Narasimhan
filtration (Section~4.2, loc.~cit.).

(ii) The wall-crossing formula of Kontsevich--Soibelman 2008
Theorem~3.1 is an equality of products in $G_{\mathrm{KS}}(X)$
across adjacent chambers; by taking the logarithm we obtain the gauge
transformation $\Theta_{\zeta'} = e^{\mathrm{ad}_{\alpha_W}}(\Theta_\zeta)$
with $\alpha_W$ the infinitesimal generator supported on the wall class
(Section~2.7, loc.~cit., Lemma~2.8).

(iii) The MC equation is the closedness of the motivic DT generating
series under the Hall product (Joyce--Song 2011
\textup{arXiv:0810.5645} Theorem~5.3 for the numerical case; KS 2008
Section~4 for the motivic case).

(iv) The pentagon is the $A_2$ quantum dilogarithm identity
(Faddeev--Kashaev, Faddeev--Volkov), which in the Hall algebra of the
$A_2$ quiver is the five-term relation
$\mathbb{E}(e_1) \mathbb{E}(e_2) = \mathbb{E}(e_2) \mathbb{E}(e_1 + e_2)
\mathbb{E}(e_1)$ where $\mathbb{E}(x) = \prod_k (1 - q^{k+1/2} x)^{-1}$
is the quantum dilogarithm (Reineke 2011 \textup{arXiv:0903.0261}
Theorem~5.1). Algebraic verification through height~$2$ is
programmed in \texttt{compute/tests/test\_conifold\_wall\_crossing.py}
(73 tests).
\end{proof}

\begin{remark}[Relation to the bar-side MC in Theorem~\ref{thm:wall-crossing-mc}]
\label{rem:two-mc-equations}
Theorem~\ref{thm:wall-crossing-mc} in Chapter~\ref{ch:toric-coha}
states the MC gauge equivalence
in $\mathfrak{g}^{\mathrm{mod}}_{A_X}$ (convolution dgLA of the chiral
algebra $A_X$). The present Theorem~\ref{thm:ks-wall-crossing-mc-on-coha-dgla}
states it in $\mathfrak{g}_{\mathrm{KS}}(X)$ (motivic Hall dgLA). These
are two MC equations on two different dg Lie algebras, linked by the
chiralisation functor $\Phi$:
\[
  \Phi_* \colon \mathfrak{g}_{\mathrm{KS}}(X)
   \longrightarrow \mathfrak{g}^{\mathrm{mod}}_{A_X}.
\]
The compatibility ``$\Phi_*(\Theta_{\zeta}^{\mathrm{KS}})
= \Theta_\zeta^{\mathrm{mod}}$'' is the CONTENT of the
CoHA-to-chiral bridge; it is PROVED for $X = \mathbb{C}^3$ (via the
Schiffmann--Vasserot identification plus the explicit Feynman-diagram
match of Section~\ref{sec:5d-cs-toric-cy3}) and PROVED for toric
CY$_3$ without compact $4$-cycles conditional on the RSYZ
identification and the $\mathbb{C}^3$-local chart-gluing of
Theorem~\ref{thm:toric-chart-gluing}. Outside this locus it is
conjectural and is part of CY-A$_3$'s content.
\end{remark}

% ========================================================================
\section{Conifold wall-crossing: cluster algebra and chiral conjecture}
\label{sec:conifold-cluster-wall-crossing}

\begin{theorem}[Conifold wall-crossing as cluster mutation]
\label{thm:conifold-cluster-wall-crossing}
\ClaimStatusProvedHere
Let $X_{\mathrm{con}}$ be the resolved conifold
$\mathcal{O}(-1) \oplus \mathcal{O}(-1) \to \mathbb{P}^1$ with
Klebanov--Witten quiver $(Q_{\mathrm{con}}, W_{\mathrm{con}})$. Then:
\begin{enumerate}[label=\textup{(\roman*)}]
  \item The motivic DT wall-crossing at the single wall
    $\gamma_0 = (1, 0) - (0, 1) \in K_0(\mathrm{Coh}(X_{\mathrm{con}}))$
    is the exponential gauge action of Theorem~%
    \ref{thm:ks-wall-crossing-mc-on-coha-dgla}(ii).
  \item This wall-crossing matches the cluster-mutation structure of
    the $A_2$ cluster algebra: the two chambers of the stability space
    correspond to the two exchange paths around the $A_2$ pentagon.
  \item The bar-complex dimensions are
    $\dim B^k_{\mathrm{I}}(Y^+) = 2^k$ (resolved chamber) and
    $\dim B^k_{\mathrm{II}}(Y^+) = 3^k$ (flopped chamber), verified
    through degree $k = 12$ (\textup{124} tests in
    \texttt{conifold\_bar\_complex.py}).
  \item The lift to an $E_1$-chiral algebra identity (i.e.,
    chiralisation of the cluster structure) is PROVED for $\mathbb{C}^3$
    local charts and conditional outside them.
\end{enumerate}
\end{theorem}

\begin{proof}
(i) is Theorem~\ref{thm:ks-wall-crossing-mc-on-coha-dgla}(ii) applied
to the conifold quiver; the single wall class
$\gamma_0 = (1, 0) - (0, 1)$ is the primitive wall in the unique
two-dimensional stability space, and the wall-crossing generator is
$\alpha_{\gamma_0}$ with Euler form $\chi(\gamma_0, \gamma_0) = 2$
(two non-canceling Ext$^1$ contributions from the superpotential).

(ii) is Nagao--Nakajima 2011 \textup{arXiv:0809.2992} Theorem~3.5:
the conifold stability space has two chambers related by a single
mutation, and the pentagon periodicity is the $A_2$ cluster algebra
mutation class of rank $2$.

(iii) by direct bar-complex enumeration (Construction~%
\ref{constr:conifold-bar} in Chapter~\ref{ch:toric-coha}, with the
engine verification cited).

(iv) conditional: the chiralisation of cluster algebras is a
conjectural content of the programme; the local $\mathbb{C}^3$ chart
lifts are proved via Schiffmann--Vasserot, but the global assembly for
the conifold is a chart-gluing conditional on the local-to-global
descent for $E_1$-chiral algebras
(Theorem~\textup{\ref{thm:toric-chart-gluing}} plus
Proposition~\textup{\ref{prop:mutation-e1-equivalence}}).
\end{proof}

% ========================================================================
\section{Positive half versus full Yangian}
\label{sec:y-plus-vs-full-y}

\begin{corollary}[$Y^+$ is not the full Yangian; Drinfeld double yields it]
\label{cor:y-plus-vs-full-Y}
\ClaimStatusProvedHere
\begin{enumerate}[label=\textup{(\roman*)}]
  \item $\mathcal{H}(\mathbb{C}^3) \simeq Y^+(\widehat{\mathfrak{gl}}_1)$
    is the POSITIVE HALF of the affine Yangian
    $Y(\widehat{\mathfrak{gl}}_1)$
    (Schiffmann--Vasserot 2013 Theorem~A).
  \item The full affine Yangian
    $Y(\widehat{\mathfrak{gl}}_1) = Y^+ \otimes Y^0 \otimes Y^-$
    is the Drinfeld double
    $D(\mathcal{H}(\mathbb{C}^3))$ with respect to the non-degenerate
    Hopf pairing of Procha\'zka--Rap\v{c}\'ak 2018
    \textup{arXiv:1807.11304}, Theorem~3.6.
  \item The identification $Y(\widehat{\mathfrak{gl}}_1)
    \simeq \mathcal{W}_{1+\infty}$ at the self-dual level
    $\psi = 1$ (Procha\'zka--Rap\v{c}\'ak) is a statement about the
    FULL double, not about the positive half alone; it makes crucial
    use of the triangular decomposition $Y = Y^+ \otimes Y^0 \otimes Y^-$.
  \item For general toric CY$_3$ $X$ without compact $4$-cycles, the
    CoHA $\mathcal{H}(Q_X, W_X) \simeq Y^+(\widehat{\mathfrak{g}}_{Q_X})$
    identifies with the positive half of an affine super Yangian (RSYZ
    2020); the full affine super Yangian is obtained as the Drinfeld
    double of $\mathcal{H}$, conditional on a non-degenerate Hopf pairing
    which is proved for $\mathbb{C}^3$ and conjectural for general toric
    quivers.
\end{enumerate}
\end{corollary}

\begin{proof}
(i) is Schiffmann--Vasserot 2013 Theorem~A, statement that the Hall
product on $\mathcal{H}(\mathbb{C}^3)$ agrees with the Drinfeld
coproduct's positive-half multiplication on
$Y^+(\widehat{\mathfrak{gl}}_1)$ specialised at generic central charge.

(ii) is Procha\'zka--Rap\v{c}\'ak 2018 Section~3. The Drinfeld double
$D(Y^+)$ is constructed from $Y^+$ and its graded dual $(Y^+)^*$ via
the Hopf pairing; the resulting algebra contains both copies as
Hopf subalgebras and has cross-relations determined by the pairing.
The identification of $(Y^+)^*$ with $Y^-(\widehat{\mathfrak{gl}}_1)$
(the CoHA of the opposite quiver) uses the CY$_3$-duality
Koszul pairing.

(iii) by direct inspection of the Procha\'zka--Rap\v{c}\'ak
isomorphism: their identification is between the full Drinfeld
double and $\mathcal{W}_{1+\infty}$, which requires both positive
and negative modes. The positive half $Y^+$ alone is isomorphic to the
positive-mode subalgebra of $\mathcal{W}_{1+\infty}$, but the full
W-algebra structure requires the double.

(iv) for general toric $X$, RSYZ 2020 \textup{arXiv:2004.07841}
Theorem~1.1 proves
$\mathcal{H}(Q_X, W_X) \simeq Y^+(\widehat{\mathfrak{g}}_{Q_X})$ with
$\widehat{\mathfrak{g}}_{Q_X}$ the affine super Lie algebra determined
by the toric fan. The Drinfeld double construction requires a
non-degenerate Hopf pairing between $Y^+$ and its dual; this is proved
for $\mathbb{C}^3$ (SV + PR) and for a handful of specific toric
examples (e.g., conifold via Klebanov--Witten), but the general case
remains conjectural. In
Theorem~\ref{thm:toric-cy3-chiral-qg}(III) the phrase ``the full
affine super Yangian'' should be read as ``the Drinfeld double of the
CoHA'', with the identification with a specific known Yangian being
PROVED for $\mathbb{C}^3$ and CONDITIONAL on a Hopf-pairing lemma
elsewhere.
\end{proof}

\begin{remark}[What the ``Drinfeld double'' gives us]
\label{rem:drinfeld-double-content}
The Drinfeld double $D(\mathcal{H})$ of an associative graded algebra
$\mathcal{H}$ with a compatible coideal-coalgebra structure (the
Yang--Zhao topological coproduct) is always defined abstractly, as
$\mathcal{H} \otimes \mathcal{H}^{\mathrm{cop}, *}$ with cross-relations
from the Hopf pairing. The content of Procha\'zka--Rap\v{c}\'ak for
$\mathbb{C}^3$ is that THIS abstract double is concretely identifiable
with $\mathcal{W}_{1+\infty}$ at $\psi = 1$. For general toric CY$_3$,
the abstract double is always available; the identification with a
specific named algebra (an affine super Yangian of a specific super
Lie algebra) is the RSYZ content and requires a Hopf pairing whose
non-degeneracy has been verified in the toric cases covered by RSYZ
2020 and left conjectural for exotic toric quivers with 4-cycle
singularities.
\end{remark}

% ========================================================================
\section{Reconciliation ledger}
\label{sec:reconciliation-ledger-coha}

This final section is the ledger that reconciles the
content of this chapter with the structural claims of
Chapter~\ref{ch:toric-coha}. Nothing in the structural assembly of
Theorem~\ref{thm:toric-cy3-chiral-qg} changes; what changes is the
precise reading of three phrases:

\begin{enumerate}[label=\textup{(\Roman*)}]
  \item \emph{``$D^2 = 0$ in the CoHA differential''} (clause (iii) of
    Theorem~\ref{thm:wall-crossing-mc}). Read as:
    $\partial_W^2 = 0$ in the Ginzburg dg algebra $\Pi(Q, W)$,
    equivalently $d_B^2 = 0$ in the bar complex
    $B^{\mathrm{ord}}(\mathcal{H})$. The CoHA $\mathcal{H}$ itself is
    an associative graded algebra without internal differential.
  \item \emph{``topological bialgebra''} (Component~(I) of
    Theorem~\ref{thm:toric-cy3-chiral-qg}). Read as:
    $\mathcal{H}(Q_X, W_X)$ has the Hall product (associative) and a
    coideal-subalgebra topological coproduct (Yang--Zhao
    localisation); the resulting structure is a coideal topological
    bialgebra, whose Drinfeld double is the full affine super
    Yangian. The word ``bialgebra'' is used in the coideal-topological
    sense, not the strict Hopf-algebra sense.
  \item \emph{``the full affine super Yangian''} (Component~(III) of
    Theorem~\ref{thm:toric-cy3-chiral-qg}). Read as: the Drinfeld
    double $D(\mathcal{H}(Q_X, W_X))$, which coincides with a specific
    named affine super Yangian for $\mathbb{C}^3$ (Procha\'zka--Rap\v{c}\'ak),
    the conifold (Klebanov--Witten / RSYZ), local $\mathbb{P}^2$
    (RSYZ), and local $\mathbb{P}^1 \times \mathbb{P}^1$ (RSYZ); for
    exotic toric quivers the identification with a named Yangian is
    conditional on a Hopf-pairing non-degeneracy lemma.
\end{enumerate}

With these three reinterpretations, the structural assembly in Chapter~%
\ref{ch:toric-coha} stands as presented. The assembled
chiral quantum group datum of Theorem~\ref{thm:toric-cy3-chiral-qg}
stands; this chapter records the precise ledger.

\begin{remark}[Cross-reference with Remark~\ref{rem:coha-bar-macmahon}]
\label{rem:cross-ref-macmahon}
The MacMahon-character remark
Remark~\ref{rem:coha-bar-macmahon} of Chapter~\ref{ch:toric-coha}
already states the algebra-vs-coalgebra distinction correctly
(``the CoHA and the bar complex remain distinct mathematical
objects (one is an algebra, the other its coalgebra resolution),
and the character coincidence is a \emph{consequence} of the isomorphism,
not an identification''). The present chapter extends this discipline
from the single MacMahon remark to the full wall-crossing treatment.
\end{remark}

\bigskip
\noindent\textit{Verification}: see
\texttt{compute/tests/test\_coha\_wall\_crossing\_platonic.py}. Four
\texttt{\textbackslash ClaimStatusProvedHere} decorators with pairwise-disjoint
\texttt{derived\_from} / \texttt{verified\_against} sources covering
Kontsevich--Soibelman 2008 (motivic DT), Schiffmann--Vasserot 2013
(CoHA$(\mathbb{C}^3) = Y^+$), Davison 2015 (positivity and integrality),
and Procha\'zka--Rap\v{c}\'ak 2018 (Drinfeld-double~$=\mathcal{W}_{1+\infty}$).
